\section{Giới thiệu}

Trong nhiều lĩnh vực khoa học và kỹ thuật hiện nay, dữ liệu thường xuất hiện
dưới dạng tín hiệu như âm thanh, hình ảnh, sóng điện từ hoặc dữ liệu cảm biến.
Một tín hiệu có thể thay đổi theo thời gian và chứa nhiều thông tin khác nhau.
Để hiểu và phân tích tín hiệu, người ta cần xác định những thành phần tạo nên
tín hiệu đó, đặc biệt là các thành phần tần số.

\subsection{Biến đổi Fourier}

Biến đổi Fourier (FT) là một trong những công cụ toán học quan trọng dùng để
phân tích tín hiệu. FT là phương pháp chuyển đổi tín hiệu từ miền thời gian
sang miền tần số.

\begin{itemize}
    \item Miền thời gian (Time Domain): biểu diễn tín hiệu theo sự thay đổi
    của nó theo thời gian.
    \item Miền tần số (Frequency Domain): biểu diễn tín hiệu theo các thành
    phần tần số tạo nên tín hiệu đó.
\end{itemize}

Ví dụ, một bài hát có thể được xem là sự kết hợp của nhiều âm thanh với các
tần số khác nhau. FT giúp tách các thành phần này để phân tích riêng biệt.

Đối với tín hiệu liên tục $x(t)$, FT được định nghĩa bởi công thức:
\begin{equation}
  X(f) = \int_{-\infty}^{\infty} x(t)e^{-i2\pi ft}\,dt,
\end{equation}
trong đó:
\begin{itemize}
    \item $x(t)$: tín hiệu trong miền thời gian.
    \item $X(f)$: tín hiệu sau khi chuyển sang miền tần số.
    \item $f$: tần số.
    \item $e^{-i2\pi ft}$: hàm mũ phức dùng để biểu diễn sóng điều hòa.
    \item $i$: đơn vị ảo với $i^2 = -1$.
\end{itemize}
Công thức trên cho phép phân tích tín hiệu thành tổng hợp của nhiều sóng sin
và cos với các tần số khác nhau.

\subsection{Biến đổi Fourier rời rạc}

Trong thực tế, máy tính và các thiết bị số không xử lý tín hiệu liên tục mà xử
lý dữ liệu rời rạc dưới dạng các mẫu (samples). Vì vậy, người ta sử dụng biến
đổi Fourier rời rạc (DFT).

DFT chuyển đổi một dãy dữ liệu rời rạc từ miền thời gian sang miền tần số. Nếu
tín hiệu có $N$ mẫu dữ liệu:

\begin{equation}
  X(k)=\sum_{n=0}^{N-1}x(n)e^{-i\frac{2\pi}{N}kn},
  \quad k=0,1,\ldots,N-1,
\end{equation}
trong đó:
\begin{itemize}
    \item $x(n)$: giá trị tín hiệu tại mẫu thứ $n$.
    \item $X(k)$: thành phần tần số thứ $k$.
    \item $N$: tổng số mẫu dữ liệu.
    \item $e^{-i\frac{2\pi}{N}kn}$: hệ số quay phức dùng để phân tích tần số.
\end{itemize}

DFT cho phép xác định các tần số xuất hiện trong tín hiệu và cường độ của từng
tần số. Đây là nền tảng quan trọng trong xử lý tín hiệu số, xử lý âm thanh và
xử lý ảnh.

Tuy nhiên, việc tính DFT trực tiếp cần rất nhiều phép tính. Nếu tín hiệu có
$N$ mẫu thì cần khoảng $N^2$ phép toán. Khi số lượng dữ liệu lớn, thời gian xử
lý sẽ tăng rất nhanh.

\subsection{Biến đổi Fourier nhanh}

Để giải quyết vấn đề tốc độ tính toán của DFT, thuật toán biến đổi Fourier
nhanh (FFT) được phát triển. FFT là thuật toán dùng để tính DFT nhanh hơn bằng
cách chia bài toán lớn thành nhiều bài toán nhỏ hơn. Nhờ đó, số lượng phép tính
được giảm đáng kể và tốc độ xử lý được cải thiện rõ rệt.

FFT không làm thay đổi kết quả của DFT mà chỉ tối ưu cách tính toán. Thay vì
thực hiện toàn bộ phép tính trực tiếp, FFT tận dụng các tính chất đối xứng và
tuần hoàn của FT để giảm số phép toán cần thiết. Độ phức tạp tính toán của DFT
trực tiếp là $O(N^2)$, trong khi FFT giảm xuống còn $O(N\log N)$.

Nhờ khả năng xử lý nhanh và hiệu quả, FFT trở thành một trong những thuật toán
quan trọng nhất trong khoa học máy tính và xử lý tín hiệu số. FFT hiện được
ứng dụng rộng rãi trong xử lý âm thanh, xử lý ảnh, truyền thông số, y học,
ra-đa và nhiều lĩnh vực kỹ thuật khác.

\subsection{Biến đổi Fourier ngược}

Sau khi tín hiệu được chuyển từ miền thời gian sang miền tần số bằng biến đổi
Fourier, ta có thể thực hiện quá trình ngược lại để khôi phục tín hiệu ban đầu.
Phép biến đổi này được gọi là biến đổi Fourier ngược (IFT).

Nếu FT giúp phân tích tín hiệu $x(t)$ thành các thành phần tần số $X(f)$, thì
IFT dùng các thành phần tần số đó để tổng hợp lại tín hiệu trong miền thời
gian. Đối với tín hiệu liên tục, IFT được định nghĩa bởi công thức:
\begin{equation}
  x(t)=\int_{-\infty}^{\infty}X(f)e^{i2\pi ft}\,df,
\end{equation}
trong đó:
\begin{itemize}
    \item $X(f)$: biểu diễn tín hiệu trong miền tần số.
    \item $x(t)$: tín hiệu sau khi được khôi phục trong miền thời gian.
    \item $f$: tần số.
    \item $t$: thời gian.
    \item $e^{i2\pi ft}$: hàm mũ phức dùng để tổng hợp lại tín hiệu từ các
    thành phần tần số.
\end{itemize}

Như vậy, FT và IFT là hai quá trình bổ sung cho nhau. FT dùng để phân tích tín
hiệu, còn IFT dùng để tái tạo lại tín hiệu ban đầu từ miền tần số.

\subsection{Biến đổi Fourier rời rạc ngược}

Trong thực tế, máy tính không xử lý tín hiệu liên tục mà xử lý các mẫu dữ liệu
rời rạc. Vì vậy, bên cạnh DFT, người ta cũng sử dụng biến đổi Fourier rời rạc
ngược (IDFT).

Nếu DFT chuyển dãy dữ liệu rời rạc $x(n)$ từ miền thời gian sang miền tần số
$X(k)$, thì IDFT thực hiện quá trình ngược lại, tức là khôi phục dãy tín hiệu
$x(n)$ từ các thành phần tần số $X(k)$. Đối với dãy tần số $X(k)$ gồm $N$ mẫu,
IDFT được định nghĩa bởi công thức:
\begin{equation}
  x(n)=\frac{1}{N}\sum_{k=0}^{N-1}X(k)e^{i\frac{2\pi}{N}kn},
  \quad n=0,1,\ldots,N-1,
\end{equation}
trong đó:
\begin{itemize}
    \item $X(k)$: thành phần tín hiệu trong miền tần số.
    \item $x(n)$: giá trị tín hiệu sau khi được khôi phục tại mẫu thứ $n$.
    \item $N$: tổng số mẫu dữ liệu.
    \item $k$: chỉ số thành phần tần số.
    \item $e^{i\frac{2\pi}{N}kn}$: hệ số mũ phức dùng để tổng hợp lại tín hiệu.
    \item $\frac{1}{N}$: hệ số chuẩn hóa giúp khôi phục đúng biên độ tín hiệu.
\end{itemize}

IDFT cho phép tái tạo lại tín hiệu rời rạc ban đầu từ các thành phần tần số đã
được phân tích bởi DFT. Nếu một tín hiệu được biến đổi bằng DFT rồi sau đó áp
dụng IDFT, ta có thể thu được lại tín hiệu ban đầu.

\subsection{Lịch sử phát triển}

Ý tưởng nền tảng của thuật toán FFT đã xuất hiện từ rất sớm trong lịch sử toán
học. Năm 1805, nhà toán học người Đức C. F. Gauss đã sử dụng một phương pháp
tương tự FFT để phục vụ việc tính toán quỹ đạo của các tiểu hành tinh Pallas
và Juno. Tuy nhiên, công trình này không được công bố rộng rãi vào thời điểm
đó nên chưa tạo được ảnh hưởng đáng kể đối với cộng đồng khoa học
\cite{heideman1984gauss}.

Trong giai đoạn tiếp theo, nhiều nhà khoa học tiếp tục nghiên cứu các phương
pháp cải tiến tính DFT nhằm giảm số lượng phép tính cần thực hiện. Đến năm
1942, G. C. Danielson và C. Lanczos đã đề xuất phương pháp phân rã phép DFT
trong nghiên cứu tinh thể học tia X. Công trình này đặt nền móng quan trọng
cho sự phát triển của thuật toán FFT hiện đại \cite{danielson1942}.

Bước ngoặt quan trọng nhất diễn ra vào năm 1965 khi J. W. Cooley và
J. W. Tukey công bố thuật toán FFT hiện đại \cite{cooley1965algorithm}. Hai
nhà khoa học đã chứng minh rằng DFT có thể được chia thành nhiều bài toán nhỏ
hơn, từ đó giảm đáng kể số lượng phép tính cần thực hiện. Thuật toán này giúp
giảm độ phức tạp tính toán từ $O(N^2)$ xuống còn $O(N \log N)$, làm tăng đáng
kể tốc độ xử lý dữ liệu.

Sự ra đời của FFT đã tạo nên bước tiến lớn trong khoa học máy tính và xử lý
tín hiệu số. Nhờ khả năng tính toán nhanh và hiệu quả, FFT nhanh chóng được
ứng dụng rộng rãi trong nhiều lĩnh vực như xử lý âm thanh, xử lý ảnh, truyền
thông số, ra-đa, y học và các ngành khoa học kỹ thuật khác.

Ngày nay, FFT được xem là một trong những thuật toán quan trọng nhất của thế
kỷ XX. Nhà toán học Gilbert Strang từng nhận định FFT là “thuật toán số quan
trọng nhất của thời đại chúng ta” \cite{strang1993wavelet}. Đồng thời, FFT
cũng được xếp vào danh sách 10 thuật toán hàng đầu trong số đặc biệt của tạp
chí \textit{IEEE Computing in Science \& Engineering}
\cite{10.1109/MCISE.2000.814652}.
